“Heterodyne” is the right word if what you had in mind was “there is a carrier sitting on top of the thing, and multiplying by the opposite phase removes it.” In the complex-demodulation formulation, one explicitly multiplies by \(e^{-i\omega_0(t-u)}\) to shift spectral content near \(\omega_0\) down to baseband, then studies the slowly varying envelope that remains. [ppl-ai-file-upload.s3.amazonaws](https://ppl-ai-file-upload.s3.amazonaws.com/web/direct-files/attachments/37918610/d8f28071-4d79-4905-9907-834b898a8661/Evolutionary-spectra-and-non-stationary-processes.pdf)

## Heterodyne meaning

In plain terms, heterodyning is frequency translation: a band centered at \(\omega_0\) is moved to \(0\) by mixing with \(e^{-i\omega_0 t}\), and moving it back up is just the inverse multiplication by \(e^{i\omega_0 t}\). The discussion of complex demodulation in the source says exactly this in operational form: form \(X_t e^{i\omega t}\) or \(X_t e^{-i\omega t}\), smooth, and the resulting local power estimates describe the moving spectrum near that carrier. [ppl-ai-file-upload.s3.amazonaws](https://ppl-ai-file-upload.s3.amazonaws.com/web/direct-files/attachments/37918610/d8f28071-4d79-4905-9907-834b898a8661/Evolutionary-spectra-and-non-stationary-processes.pdf)

That is why the extra exponential matters so much. A covariance or kernel with a factor like \(e^{-i\Delta}\) is not “plain sinc at baseband”; it is what baseband looks like before demodulation, with a residual carrier phase still attached. [ppl-ai-file-upload.s3.amazonaws](https://ppl-ai-file-upload.s3.amazonaws.com/web/direct-files/attachments/37918610/d8f28071-4d79-4905-9907-834b898a8661/Evolutionary-spectra-and-non-stationary-processes.pdf)

## Why it fits

If you start from a baseband process with sinc-type covariance and then multiply the process by a pure phase \(e^{it}\), the second-order structure picks up exactly a phase factor depending on the lag, because modulation in time is frequency shift in the spectrum and demodulation removes it again. That is the whole heterodyne connection: the exponential is the carrier bookkeeping term, not an incidental decoration. [ppl-ai-file-upload.s3.amazonaws](https://ppl-ai-file-upload.s3.amazonaws.com/web/direct-files/attachments/37918610/d8f28071-4d79-4905-9907-834b898a8661/Evolutionary-spectra-and-non-stationary-processes.pdf)

So the reason that word probably stuck in your head is that the structure is not “just a bandlimited kernel”; it is “a bandlimited kernel seen in a rotated frame.” Once you say that, the exponential stops being an annoyance and becomes the diagnostic that tells you a demodulation step is available. [ppl-ai-file-upload.s3.amazonaws](https://ppl-ai-file-upload.s3.amazonaws.com/web/direct-files/attachments/37918610/d8f28071-4d79-4905-9907-834b898a8661/Evolutionary-spectra-and-non-stationary-processes.pdf)

## If you frame \(Z\)

If you frame \(Z\) as an oscillatory process, then the heterodyne picture says you should look for a representation as “slow envelope times fast phase,” because the source’s oscillatory-process form is exactly \(A_t(\omega)e^{i\omega t}\) with the slowly varying part carried by the envelope and the fast oscillation carried by the exponential. In that reading, demodulation is the move that strips off the fast phase so you can inspect the underlying baseband object directly. [ppl-ai-file-upload.s3.amazonaws](https://ppl-ai-file-upload.s3.amazonaws.com/web/direct-files/attachments/37918610/d8f28071-4d79-4905-9907-834b898a8661/Evolutionary-spectra-and-non-stationary-processes.pdf)

That is also why the exact Riemann-Siegel-type decomposition you were pushing toward feels heterodyne-adjacent: discrete oscillatory terms plus a continuous remainder already look like spectral content written in a carrier-bearing form, and the natural next move is to ask what survives after the deterministic phase is peeled off. [ppl-ai-file-upload.s3.amazonaws](https://ppl-ai-file-upload.s3.amazonaws.com/web/direct-files/collection_5daedf3c-ee31-4d76-80d9-57772ecae8fc/53d6aa77-5d05-4969-96c3-a5cd09fdfe13/check-out-this-sustained-rise-MsM6kVIXRh2siTVmSiMI9A.md)
